\chapter{DESIGN AND CONSTRUCTION}

\section{Mechanical Design}

\subsection{Ballast Tank System}

The ballast tank (Figure~\ref{fig:1.1}) utilizes a 500ml syringe for 
ease of manufacturing and control. An M6 lead screw guides 
a threaded coupling along the rod, which is coupled to a 
JGB37-520 motor at the terminal. The ballast tank is divided 
into two main sections: the mobile section and the stationary 
section.

\vspace{1em}

\begin{figure}[H]
\begin{center}
    \includegraphics[width=6cm,height=3.375cm]{BallastMov.png}
\end{center}

\caption{Mobile Section of Ballast Tank}
\label{fig:3.1}
\end{figure}

\vspace{1em}

The mobile section (Figure~\ref{fig:3.1}) consists of two 3D-printed 
parts connected by three M4 lead screws, each 100mm in length. 
The front-facing part of the syringe has a rubber seal, while 
the aft section features a threaded coupling. This section can 
move independently along the syringe’s radial axis.

\vspace{1em}

\begin{figure}[H]
\begin{center}
    \includegraphics[width=6cm,height=3.375cm]{BallastSta.png}
\end{center}

\caption{Stationary Section of Ballast Tank}
\label{fig:3.2}
\end{figure}

\vspace{1em}

The stationary section (Figure~\ref{fig:3.2}) is the main syringe 
body, with its aft end joined to a 3D-printed cap with a snap-fit 
mechanism. This cap connects to the motor and houses the motor-M6 
coupler. Three M3 lead screws are bonded to the end cap, acting as 
guiding rails for the mobile section. At the tip of each screw, a 
3D-printed cap holds them perpendicular, and a bearing is attached 
to ensure smooth rotation of the main M6 lead screw.

\subsection{Motor Tube Compartment}

Figure~\ref{fig:1.2}, the propulsion system utilizes A2212 930kV 
brushless DC (BLDC) motors, each securely mounted onto a coated 
structural component positioned between two acrylic tubes. These 
tubes have an outer diameter (OD) of 50 mm, a thickness of 3 mm, 
and a length of 70 mm. At the end of each tube, an outlet 
facilitates controlled water intake and expulsion, enabling 
efficient fluid flow management within the system.

\vspace{1em}

\begin{figure}[H]
\begin{center}
    \includegraphics[width=6cm,height=3.375cm]{Prop.png}
\end{center}

\caption{Propellor Design}
\label{fig:3.3}
\end{figure}

\vspace{1em}

Figure~\ref{fig:3.3} illustrates the propeller design, which 
is a 3D-printed part coated with resin layer. The propeller 
is designed to have 8 blades. The blades are angled to optimize 
thrust while minimizing drag, ensuring efficient propulsion 
for the underwater vehicle.

The propeller also have a circular ring around the blades, 
which serves to enhance the structural integrity of the propeller 
. The ring helps to stabilize the flow of water around the blades, 
creating a streamline laminar output flow and reducing turbulence, 
which can improve the overall noise performance of the propulsion 
system.

\subsection{Circuit Compartment}

The circuit tube is designed to be partially sealed, leaving only 
the necessary space for cable connections. This design ensures a 
partial safety precaution against small leaks, with all core 
components securely fitted inside to prevent interference with 
other parts of the robot.

\vspace{1em}

\begin{figure}[H]
\begin{center}
    \includegraphics[width=4.8cm,height=3.7cm]{Circuit.png}
\end{center}

\caption{Circuit Compartment Design}
\label{fig:3.4}
\end{figure}

\vspace{1em}

Figure~\ref{fig:3.4} shows the circuit compartment design. 
The circuit is housed within an acrylic tube with an outer diameter 
(OD) of 90 mm, a wall thickness of 3 mm, and a length of 220 mm. At 
each end of the tube, a 3D-printed cap is securely snapped into 
place, enclosing the compartment while also serving as a connector 
for a 100 mm-long 1010 aluminum profile. These aluminum profiles 
function as mounting structures for the sensor system, providing 
stability and ease of integration.

Enclosed within the acrylic tube is the primary control and power 
management system, consisting of a Raspberry Pi, an Arduino Nano, 
voltage sensors, an inertial measurement unit (IMU), a terminal, 
and a Lipo Battery for rechargeability and stability. These 
components collectively form the core circuitry housed within 
the enclosure.

\vspace{1em}

\begin{figure}[H]
\begin{center}
    \includegraphics[width=4.8cm,height=3.7cm]{CircuitSide.png}
\end{center}

\caption{Circuit Compartment Design Side View}
\label{fig:3.5}
\end{figure}

\vspace{1em}

To facilitate maintenance, all components are mounted on an lead 
3D-printed rod, enabling a controlled sliding motion for easy access 
and adjustments.

Externally, the system integrates a laser distance sensor, hydrostatic 
pressure sensors, and a Raspberry Pi camera, enabling environmental 
perception and data acquisition for autonomous operation.

\subsection{Main Frame Design}

\vspace{1em}

\begin{figure}[H]
\begin{center}
    \includegraphics[width=4.8cm,height=3.7cm]{Frame.png}
\end{center}

\caption{Main Frame Design}
\label{fig:3.6}
\end{figure}

\vspace{1em}

The robot's shell is divided into four distinct sections, each 
separated by an acrylic plane: the front dome, which encapsulates 
the Raspberry Pi Camera for visual data acquisition and the 
fail-safe module to ensure system redundancy; two central 
compartments designed to house the ballast tank, propulsion 
motor, and circuit enclosure; and the aft dome.

This modular design, with acrylic planes separating the 
compartments, optimizes component placement, provides easy 
access for maintenance and system integration, and ensures 
structural integrity and functionality. The central coupler is 
coated and bonded to two acrylic tubes with an outer diameter 
(OD) of 220 mm, a wall thickness of 3 mm, and a length of 10 mm. 
At the end of each tube, an additional coated coupler connects 
the dome to the tube, ensuring structural integrity.

Four M3 lead screws run along the sides of the robot, providing 
compression to maintain structural stability. The dome is designed 
for removal, as it is not permanently sealed. Instead, an O-ring 
(or a Rubber Seal), in combination with the applied compression, 
ensures a watertight seal while allowing for easy access to 
internal components.

\section{Electrical Design}
\section{Control System Design}
\section{AI Detection System Design}
\section{Manufacturing Process}
